% Methodology
\section{Methodology}
\label{sec:methodology}

Building on our hypothesis that truthfulness benchmarks measure distinct capability dimensions, we design a correlation analysis framework to quantify the relationship structure between TruthfulQA, HaluEval, and FactScore across a diverse model population.

\subsection{Overview}

Our methodology comprises four components: (1) model population selection ensuring diversity for robust correlation estimation, (2) benchmark score collection using standardized evaluation, (3) correlation analysis with appropriate statistical methods, and (4) factor analysis to test multi-dimensional structure. We structure our analysis around four hypotheses testing progressively specific claims about the independence of truthfulness dimensions.

\subsection{Model Population}

\textbf{Rationale:} Correlation estimates require sufficient sample size and population diversity. A homogeneous population (e.g., only 7B Llama variants) would yield unstable estimates and limit generalizability.

We select N=50 models from the Open LLM Leaderboard satisfying the following criteria:

\begin{table}[h]
\centering
\begin{tabular}{@{}lll@{}}
\toprule
\textbf{Criterion} & \textbf{Values} & \textbf{Rationale} \\
\midrule
Architecture & Llama, Mistral, Falcon, Phi, Qwen, Gemma, Yi & Cover major decoder-only families \\
Scale & 7B, 13B, 34B, 70B & Test scale effects on correlation structure \\
Training variant & Base, Instruct, Chat, DPO & Include fine-tuning diversity \\
Availability & Public weights on HuggingFace & Reproducibility \\
\bottomrule
\end{tabular}
\end{table}

This selection yields 7 architectures $\times$ 4 scales $\times$ 4 variants, with sampling to reach N=50 where all combinations are not available. The population intentionally excludes proprietary models (GPT-4, Claude) due to lack of standardized benchmark access.

\subsection{Benchmark Score Collection}

\textbf{Rationale:} Correlation analysis requires scores from the same models on all benchmarks under consistent evaluation conditions.

We collect scores for:

\begin{itemize}
\item \textbf{TruthfulQA-MC2:} Multiple-choice accuracy on 817 questions across 38 categories. MC2 allows multiple correct answers, providing finer-grained measurement than MC1.
\item \textbf{HaluEval:} Binary detection accuracy across QA, dialogue, and summarization subtasks (35,000 samples total).
\item \textbf{FactScore:} Atomic fact precision on biography generation task.
\item \textbf{MMLU:} General knowledge baseline (57 subjects, 14,000 questions) for comparison.
\end{itemize}

For TruthfulQA and MMLU, we use lm-evaluation-harness with 5-shot prompting. For HaluEval, we use the official evaluation code. For FactScore, we use proxy methodology validated against available official scores due to computational constraints of full atomic decomposition.

\subsection{Correlation Analysis}

\textbf{Rationale:} We choose Spearman correlation over Pearson because benchmark scores may have non-normal distributions and we want rank-based relationships robust to outliers.

\subsubsection{Primary Analysis}

For each benchmark pair, we compute:

\begin{enumerate}
\item \textbf{Spearman correlation coefficient ($\rho$):} Measures monotonic relationship between benchmark rankings.
\item \textbf{95\% confidence interval:} Via 1,000 bootstrap iterations.
\item \textbf{Statistical significance:} p-values with Bonferroni correction for multiple comparisons ($\alpha = 0.05/6 = 0.0083$ for 6 pairwise comparisons).
\end{enumerate}

\subsubsection{Threshold Interpretation}

We define correlation ranges:
\begin{itemize}
\item \textbf{r $<$ 0.3:} Low correlation (approaching independence)
\item \textbf{0.3 $\leq$ r $<$ 0.7:} Moderate correlation (partial independence)
\item \textbf{r $\geq$ 0.7:} High correlation (approaching interchangeability)
\end{itemize}

The critical test: inter-benchmark correlations should exceed the unrelated-benchmark baseline (r(MMLU-Physics, HaluEval) $\approx$ 0.10) while remaining below 0.7.

\subsubsection{Baseline Reference}

We establish a baseline correlation using benchmarks measuring clearly distinct constructs:
\begin{itemize}
\item r(MMLU-Physics, HaluEval) serves as the ``unrelated benchmark'' floor.
\item Correlations between truthfulness benchmarks should exceed this baseline (otherwise they might be measuring unrelated constructs) but fall below intra-benchmark correlations (otherwise they would be interchangeable).
\end{itemize}

\subsection{Factor Analysis}

\textbf{Rationale:} Correlation matrices reveal pairwise relationships; factor analysis reveals latent structure. If truthfulness is a single construct, one factor should explain $>$80\% of variance. If multi-dimensional, 2--3 factors will be needed.

We apply Principal Component Analysis (PCA) with:
\begin{itemize}
\item \textbf{Inputs:} Standardized benchmark scores (z-scores)
\item \textbf{Components:} Extract components with eigenvalue $>$ 1 (Kaiser criterion)
\item \textbf{Variance threshold:} Count components needed for 80\% cumulative variance
\item \textbf{Interpretation:} PC1 explaining $>$80\% variance supports single-factor model; 2--3 components needed supports multi-dimensional model
\end{itemize}

\subsection{Hypothesis Structure}

We organize our analysis around four testable hypotheses with defined success criteria and gate conditions:

\begin{table}[h]
\centering
\begin{tabular}{@{}llll@{}}
\toprule
\textbf{Hypothesis} & \textbf{Type} & \textbf{Gate} & \textbf{Success Criterion} \\
\midrule
H-E1 & Existence & MUST\_WORK & $0.10 < r(\text{inter-benchmark}) < 0.7$ \\
H-M1 & Mechanism & MUST\_WORK & $r(\text{TQA, MMLU}) < r(\text{MMLU internal})$ \\
H-M2 & Mechanism & SHOULD\_WORK & $r(\text{HE, TQA}) < 0.7$ \\
H-M3 & Mechanism & SHOULD\_WORK & $r(\text{FS, TQA}) < 0.7$ AND $r(\text{FS, HE}) < 0.7$ \\
\bottomrule
\end{tabular}
\end{table}

The hypotheses form a dependency chain: H-E1 must pass before testing mechanism hypotheses. MUST\_WORK gates require passing for the overall hypothesis to be supported; SHOULD\_WORK gates allow continuation with caveats if failed.

\subsection{Implementation}

Analysis code uses:
\begin{itemize}
\item \textbf{scipy.stats.spearmanr:} Correlation computation
\item \textbf{sklearn.decomposition.PCA:} Factor analysis
\item \textbf{statsmodels:} Bootstrap confidence intervals
\end{itemize}

All code, data, and analysis scripts will be released for reproducibility.
